\documentclass[11pt,letterpaper]{article}
\newcommand{\labid}{1-B}
\newcommand{\labtitleshort}{STM32 Dev Environment}
% Shared preamble for Module 1 lab handouts.
% Usage: each handout defines \labid and \labtitleshort, then % Shared preamble for Module 1 lab handouts.
% Usage: each handout defines \labid and \labtitleshort, then % Shared preamble for Module 1 lab handouts.
% Usage: each handout defines \labid and \labtitleshort, then \input{LabHandout_Preamble}
% BEFORE \begin{document}. Compiles with a single pdflatex pass (no glossaries tool).

\usepackage[T1]{fontenc}
\usepackage[margin=1in]{geometry}
\usepackage{titlesec}
\usepackage{enumitem}
\usepackage{booktabs}
\usepackage{tabularx}
\usepackage{graphicx}
\graphicspath{{Images/}{../Images/}}
\usepackage{xcolor}
\usepackage{hyperref}
\usepackage{fancyhdr}
\usepackage{amsmath}
\usepackage{amssymb}
\usepackage{tcolorbox}
\usepackage{caption}
\tcbuselibrary{skins,breakable}

% ── Glossary (per-document, built from shared glossary.tex) ──
% Uses \makenoidxglossaries so no external makeglossaries tool is
% needed: sorting and printing is done by LaTeX itself across two
% pdflatex passes (same cadence as the TOC). Only acronyms that
% are actually \gls{}-used in the handout appear in the printed
% glossary, and each \gls{} reference is a live hyperlink that
% jumps to its entry in the Acronyms section at the end.
\usepackage[acronym,nomain,nopostdot]{glossaries}
\makenoidxglossaries
\input{../../glossary}
% Call \printhandoutglossary just before \end{document} in each handout.
\newcommand{\printhandoutglossary}{%
	\clearpage
	\printnoidxglossary[type=\acronymtype,title={Glossary}]%
}

% ── Colors (match Module 1 lesson plan) ──────────────────────
\definecolor{stevens}{RGB}{163,31,52}
\definecolor{darkgray}{RGB}{60,60,60}
\definecolor{lightgray}{RGB}{240,240,240}
\definecolor{accentblue}{RGB}{30,80,140}

\hypersetup{
	colorlinks=true,
	linkcolor=stevens,
	urlcolor=accentblue,
	citecolor=darkgray,
}

% ── Defaults for header variables (handout overrides these) ──
\providecommand{\labid}{?}
\providecommand{\labtitleshort}{Lab Handout}

% ── Running header / footer ──────────────────────────────────
\pagestyle{fancy}
\fancyhf{}
\fancyhead[L]{\small\textcolor{darkgray}{Module 1 --- Lab Handout \labid}}
\fancyhead[R]{\small\textcolor{darkgray}{\labtitleshort}}
\fancyfoot[C]{\small\thepage}
\renewcommand{\headrulewidth}{0.4pt}

% ── Section formatting ───────────────────────────────────────
\titleformat{\section}{\Large\bfseries\color{stevens}}{\thesection}{0.6em}{}
\titleformat{\subsection}{\large\bfseries\color{darkgray}}{\thesubsection}{0.6em}{}

% ── Box styles ───────────────────────────────────────────────
\tcbset{
	infobox/.style={
		enhanced, colback=blue!3, colframe=accentblue, fonttitle=\bfseries,
		title={#1}, coltitle=white, breakable, left=6pt, right=6pt, top=4pt, bottom=4pt
	},
	warning/.style={
		enhanced, colback=red!3, colframe=stevens, fonttitle=\bfseries,
		title={#1}, coltitle=white, breakable, left=6pt, right=6pt, top=4pt, bottom=4pt
	},
	deliverable/.style={
		enhanced, colback=gray!5, colframe=darkgray, fonttitle=\bfseries,
		title={#1}, coltitle=white, breakable, left=6pt, right=6pt, top=4pt, bottom=4pt
	},
}

% ── Figure helper: always draws a placeholder box (images suppressed) ─
% Usage: \labfigure{filename}{width}{caption}{label}
\newcommand{\labfigure}[4]{%
	\begin{figure}[h!]
		\centering
		\fbox{\begin{minipage}[c][2.2in][c]{#2}
			\centering\color{darkgray}\small
				\textbf{[Figure placeholder]}\\[4pt]
				\texttt{\detokenize{#1}}\\[4pt]
				\textit{#3}
		\end{minipage}}%
		\caption{#3}
		\label{#4}
	\end{figure}%
}

% ── Title banner (call at top of document body) ──────────────
\newcommand{\labheader}[2]{%
	\begin{center}
		{\LARGE\bfseries\color{stevens} Lab Handout #1}\\[3pt]
		{\Large\bfseries #2}\\[6pt]
		{\small\color{darkgray} Capstone --- Module 1: \RF{} Hardware Foundations}
	\end{center}
	\vspace{4pt}
	{\color{darkgray}\rule{\textwidth}{0.4pt}}
	\vspace{10pt}
}

% BEFORE \begin{document}. Compiles with a single pdflatex pass (no glossaries tool).

\usepackage[T1]{fontenc}
\usepackage[margin=1in]{geometry}
\usepackage{titlesec}
\usepackage{enumitem}
\usepackage{booktabs}
\usepackage{tabularx}
\usepackage{graphicx}
\graphicspath{{Images/}{../Images/}}
\usepackage{xcolor}
\usepackage{hyperref}
\usepackage{fancyhdr}
\usepackage{amsmath}
\usepackage{amssymb}
\usepackage{tcolorbox}
\usepackage{caption}
\tcbuselibrary{skins,breakable}

% ── Glossary (per-document, built from shared glossary.tex) ──
% Uses \makenoidxglossaries so no external makeglossaries tool is
% needed: sorting and printing is done by LaTeX itself across two
% pdflatex passes (same cadence as the TOC). Only acronyms that
% are actually \gls{}-used in the handout appear in the printed
% glossary, and each \gls{} reference is a live hyperlink that
% jumps to its entry in the Acronyms section at the end.
\usepackage[acronym,nomain,nopostdot]{glossaries}
\makenoidxglossaries
% Shared acronym macro container for EE800/820 lesson plan modules.
% Usage: type \IDE, \FPGA, \UART, etc. First use is expanded; later uses show acronym only.
% Requires in each document preamble:
%   \usepackage[acronym,toc]{glossaries}
%   \makeglossaries
%   \input{../glossary}
% Print used acronyms near end of document:
%   \printglossary[type=\acronymtype,title={Acronyms Used}]

% Acronym definitions (only used entries appear in the printed glossary).
\newacronym{ask}{ASK}{Amplitude Shift Keying}
\newacronym{bga}{BGA}{Ball Grid Array}
\newacronym{bom}{BOM}{Bill of Materials}
\newacronym{bram}{BRAM}{Block Random Access Memory}
\newacronym{crc}{CRC}{Cyclic Redundancy Check}
\newacronym{css}{CSS}{Chirp Spread Spectrum}
\newacronym{dma}{DMA}{Direct Memory Access}
\newacronym{dsp}{DSP}{Digital Signal Processing}
\newacronym{eda}{EDA}{Exploratory Data Analysis}
\newacronym{fifo}{FIFO}{First In, First Out}
\newacronym{fpga}{FPGA}{Field-Programmable Gate Array}
\newacronym{fsk}{FSK}{Frequency Shift Keying}
\newacronym{fsm}{FSM}{Finite State Machine}
\newacronym{gpio}{GPIO}{General-Purpose Input/Output}
\newacronym{hal}{HAL}{Hardware Abstraction Layer}
\newacronym{i2c}{I\textsuperscript{2}C}{Inter-Integrated Circuit}
\newacronym{ide}{IDE}{Integrated Development Environment}
\newacronym{ila}{ILA}{Integrated Logic Analyzer}
\newacronym{ism}{ISM}{Industrial, Scientific, and Medical}
\newacronym{isr}{ISR}{Interrupt Service Routine}
\newacronym{lut}{LUT}{Look-Up Table}
\newacronym{mcu}{MCU}{Microcontroller Unit}
\newacronym{ml}{ML}{Machine Learning}
\newacronym{ook}{OOK}{On-Off Keying}
\newacronym{per}{PER}{Packet Error Rate}
\newacronym{rf}{RF}{Radio Frequency}
\newacronym{rssi}{RSSI}{Received Signal Strength Indicator}
\newacronym{rtl}{RTL}{Register-Transfer Level}
\newacronym{snr}{SNR}{Signal-to-Noise Ratio}
\newacronym{spi}{SPI}{Serial Peripheral Interface}
\newacronym{uart}{UART}{Universal Asynchronous Receiver/Transmitter}
\newacronym{vhdl}{VHDL}{VHSIC Hardware Description Language}

% Convenience wrappers so you can type \IDE, \RF, etc.
\providecommand{\ASK}{\gls{ask}}
\providecommand{\BGA}{\gls{bga}}
\providecommand{\BOM}{\gls{bom}}
\providecommand{\BRAM}{\gls{bram}}
\providecommand{\CRC}{\gls{crc}}
\providecommand{\CSS}{\gls{css}}
\providecommand{\DMA}{\gls{dma}}
\providecommand{\DSP}{\gls{dsp}}
\providecommand{\EDA}{\gls{eda}}
\providecommand{\FIFO}{\gls{fifo}}
\providecommand{\FPGA}{\gls{fpga}}
\providecommand{\FSK}{\gls{fsk}}
\providecommand{\FSM}{\gls{fsm}}
\providecommand{\GPIO}{\gls{gpio}}
\providecommand{\HAL}{\gls{hal}}
\providecommand{\IIC}{\gls{i2c}}
\providecommand{\IDE}{\gls{ide}}
\providecommand{\ILA}{\gls{ila}}
\providecommand{\ISM}{\gls{ism}}
\providecommand{\ISR}{\gls{isr}}
\providecommand{\LUT}{\gls{lut}}
\providecommand{\MCU}{\gls{mcu}}
\providecommand{\ML}{\gls{ml}}
\providecommand{\OOK}{\gls{ook}}
\providecommand{\PER}{\gls{per}}
\providecommand{\RF}{\gls{rf}}
\providecommand{\RSSI}{\gls{rssi}}
\providecommand{\RTL}{\gls{rtl}}
\providecommand{\SNR}{\gls{snr}}
\providecommand{\SPI}{\gls{spi}}
\providecommand{\UART}{\gls{uart}}
\providecommand{\VHDL}{\gls{vhdl}}


% Call \printhandoutglossary just before \end{document} in each handout.
\newcommand{\printhandoutglossary}{%
	\clearpage
	\printnoidxglossary[type=\acronymtype,title={Glossary}]%
}

% ── Colors (match Module 1 lesson plan) ──────────────────────
\definecolor{stevens}{RGB}{163,31,52}
\definecolor{darkgray}{RGB}{60,60,60}
\definecolor{lightgray}{RGB}{240,240,240}
\definecolor{accentblue}{RGB}{30,80,140}

\hypersetup{
	colorlinks=true,
	linkcolor=stevens,
	urlcolor=accentblue,
	citecolor=darkgray,
}

% ── Defaults for header variables (handout overrides these) ──
\providecommand{\labid}{?}
\providecommand{\labtitleshort}{Lab Handout}

% ── Running header / footer ──────────────────────────────────
\pagestyle{fancy}
\fancyhf{}
\fancyhead[L]{\small\textcolor{darkgray}{Module 1 --- Lab Handout \labid}}
\fancyhead[R]{\small\textcolor{darkgray}{\labtitleshort}}
\fancyfoot[C]{\small\thepage}
\renewcommand{\headrulewidth}{0.4pt}

% ── Section formatting ───────────────────────────────────────
\titleformat{\section}{\Large\bfseries\color{stevens}}{\thesection}{0.6em}{}
\titleformat{\subsection}{\large\bfseries\color{darkgray}}{\thesubsection}{0.6em}{}

% ── Box styles ───────────────────────────────────────────────
\tcbset{
	infobox/.style={
		enhanced, colback=blue!3, colframe=accentblue, fonttitle=\bfseries,
		title={#1}, coltitle=white, breakable, left=6pt, right=6pt, top=4pt, bottom=4pt
	},
	warning/.style={
		enhanced, colback=red!3, colframe=stevens, fonttitle=\bfseries,
		title={#1}, coltitle=white, breakable, left=6pt, right=6pt, top=4pt, bottom=4pt
	},
	deliverable/.style={
		enhanced, colback=gray!5, colframe=darkgray, fonttitle=\bfseries,
		title={#1}, coltitle=white, breakable, left=6pt, right=6pt, top=4pt, bottom=4pt
	},
}

% ── Figure helper: always draws a placeholder box (images suppressed) ─
% Usage: \labfigure{filename}{width}{caption}{label}
\newcommand{\labfigure}[4]{%
	\begin{figure}[h!]
		\centering
		\fbox{\begin{minipage}[c][2.2in][c]{#2}
			\centering\color{darkgray}\small
				\textbf{[Figure placeholder]}\\[4pt]
				\texttt{\detokenize{#1}}\\[4pt]
				\textit{#3}
		\end{minipage}}%
		\caption{#3}
		\label{#4}
	\end{figure}%
}

% ── Title banner (call at top of document body) ──────────────
\newcommand{\labheader}[2]{%
	\begin{center}
		{\LARGE\bfseries\color{stevens} Lab Handout #1}\\[3pt]
		{\Large\bfseries #2}\\[6pt]
		{\small\color{darkgray} Capstone --- Module 1: \RF{} Hardware Foundations}
	\end{center}
	\vspace{4pt}
	{\color{darkgray}\rule{\textwidth}{0.4pt}}
	\vspace{10pt}
}

% BEFORE \begin{document}. Compiles with a single pdflatex pass (no glossaries tool).

\usepackage[T1]{fontenc}
\usepackage[margin=1in]{geometry}
\usepackage{titlesec}
\usepackage{enumitem}
\usepackage{booktabs}
\usepackage{tabularx}
\usepackage{graphicx}
\graphicspath{{Images/}{../Images/}}
\usepackage{xcolor}
\usepackage{hyperref}
\usepackage{fancyhdr}
\usepackage{amsmath}
\usepackage{amssymb}
\usepackage{tcolorbox}
\usepackage{caption}
\tcbuselibrary{skins,breakable}

% ── Glossary (per-document, built from shared glossary.tex) ──
% Uses \makenoidxglossaries so no external makeglossaries tool is
% needed: sorting and printing is done by LaTeX itself across two
% pdflatex passes (same cadence as the TOC). Only acronyms that
% are actually \gls{}-used in the handout appear in the printed
% glossary, and each \gls{} reference is a live hyperlink that
% jumps to its entry in the Acronyms section at the end.
\usepackage[acronym,nomain,nopostdot]{glossaries}
\makenoidxglossaries
% Shared acronym macro container for EE800/820 lesson plan modules.
% Usage: type \IDE, \FPGA, \UART, etc. First use is expanded; later uses show acronym only.
% Requires in each document preamble:
%   \usepackage[acronym,toc]{glossaries}
%   \makeglossaries
%   % Shared acronym macro container for EE800/820 lesson plan modules.
% Usage: type \IDE, \FPGA, \UART, etc. First use is expanded; later uses show acronym only.
% Requires in each document preamble:
%   \usepackage[acronym,toc]{glossaries}
%   \makeglossaries
%   \input{../glossary}
% Print used acronyms near end of document:
%   \printglossary[type=\acronymtype,title={Acronyms Used}]

% Acronym definitions (only used entries appear in the printed glossary).
\newacronym{ask}{ASK}{Amplitude Shift Keying}
\newacronym{bga}{BGA}{Ball Grid Array}
\newacronym{bom}{BOM}{Bill of Materials}
\newacronym{bram}{BRAM}{Block Random Access Memory}
\newacronym{crc}{CRC}{Cyclic Redundancy Check}
\newacronym{css}{CSS}{Chirp Spread Spectrum}
\newacronym{dma}{DMA}{Direct Memory Access}
\newacronym{dsp}{DSP}{Digital Signal Processing}
\newacronym{eda}{EDA}{Exploratory Data Analysis}
\newacronym{fifo}{FIFO}{First In, First Out}
\newacronym{fpga}{FPGA}{Field-Programmable Gate Array}
\newacronym{fsk}{FSK}{Frequency Shift Keying}
\newacronym{fsm}{FSM}{Finite State Machine}
\newacronym{gpio}{GPIO}{General-Purpose Input/Output}
\newacronym{hal}{HAL}{Hardware Abstraction Layer}
\newacronym{i2c}{I\textsuperscript{2}C}{Inter-Integrated Circuit}
\newacronym{ide}{IDE}{Integrated Development Environment}
\newacronym{ila}{ILA}{Integrated Logic Analyzer}
\newacronym{ism}{ISM}{Industrial, Scientific, and Medical}
\newacronym{isr}{ISR}{Interrupt Service Routine}
\newacronym{lut}{LUT}{Look-Up Table}
\newacronym{mcu}{MCU}{Microcontroller Unit}
\newacronym{ml}{ML}{Machine Learning}
\newacronym{ook}{OOK}{On-Off Keying}
\newacronym{per}{PER}{Packet Error Rate}
\newacronym{rf}{RF}{Radio Frequency}
\newacronym{rssi}{RSSI}{Received Signal Strength Indicator}
\newacronym{rtl}{RTL}{Register-Transfer Level}
\newacronym{snr}{SNR}{Signal-to-Noise Ratio}
\newacronym{spi}{SPI}{Serial Peripheral Interface}
\newacronym{uart}{UART}{Universal Asynchronous Receiver/Transmitter}
\newacronym{vhdl}{VHDL}{VHSIC Hardware Description Language}

% Convenience wrappers so you can type \IDE, \RF, etc.
\providecommand{\ASK}{\gls{ask}}
\providecommand{\BGA}{\gls{bga}}
\providecommand{\BOM}{\gls{bom}}
\providecommand{\BRAM}{\gls{bram}}
\providecommand{\CRC}{\gls{crc}}
\providecommand{\CSS}{\gls{css}}
\providecommand{\DMA}{\gls{dma}}
\providecommand{\DSP}{\gls{dsp}}
\providecommand{\EDA}{\gls{eda}}
\providecommand{\FIFO}{\gls{fifo}}
\providecommand{\FPGA}{\gls{fpga}}
\providecommand{\FSK}{\gls{fsk}}
\providecommand{\FSM}{\gls{fsm}}
\providecommand{\GPIO}{\gls{gpio}}
\providecommand{\HAL}{\gls{hal}}
\providecommand{\IIC}{\gls{i2c}}
\providecommand{\IDE}{\gls{ide}}
\providecommand{\ILA}{\gls{ila}}
\providecommand{\ISM}{\gls{ism}}
\providecommand{\ISR}{\gls{isr}}
\providecommand{\LUT}{\gls{lut}}
\providecommand{\MCU}{\gls{mcu}}
\providecommand{\ML}{\gls{ml}}
\providecommand{\OOK}{\gls{ook}}
\providecommand{\PER}{\gls{per}}
\providecommand{\RF}{\gls{rf}}
\providecommand{\RSSI}{\gls{rssi}}
\providecommand{\RTL}{\gls{rtl}}
\providecommand{\SNR}{\gls{snr}}
\providecommand{\SPI}{\gls{spi}}
\providecommand{\UART}{\gls{uart}}
\providecommand{\VHDL}{\gls{vhdl}}


% Print used acronyms near end of document:
%   \printglossary[type=\acronymtype,title={Acronyms Used}]

% Acronym definitions (only used entries appear in the printed glossary).
\newacronym{ask}{ASK}{Amplitude Shift Keying}
\newacronym{bga}{BGA}{Ball Grid Array}
\newacronym{bom}{BOM}{Bill of Materials}
\newacronym{bram}{BRAM}{Block Random Access Memory}
\newacronym{crc}{CRC}{Cyclic Redundancy Check}
\newacronym{css}{CSS}{Chirp Spread Spectrum}
\newacronym{dma}{DMA}{Direct Memory Access}
\newacronym{dsp}{DSP}{Digital Signal Processing}
\newacronym{eda}{EDA}{Exploratory Data Analysis}
\newacronym{fifo}{FIFO}{First In, First Out}
\newacronym{fpga}{FPGA}{Field-Programmable Gate Array}
\newacronym{fsk}{FSK}{Frequency Shift Keying}
\newacronym{fsm}{FSM}{Finite State Machine}
\newacronym{gpio}{GPIO}{General-Purpose Input/Output}
\newacronym{hal}{HAL}{Hardware Abstraction Layer}
\newacronym{i2c}{I\textsuperscript{2}C}{Inter-Integrated Circuit}
\newacronym{ide}{IDE}{Integrated Development Environment}
\newacronym{ila}{ILA}{Integrated Logic Analyzer}
\newacronym{ism}{ISM}{Industrial, Scientific, and Medical}
\newacronym{isr}{ISR}{Interrupt Service Routine}
\newacronym{lut}{LUT}{Look-Up Table}
\newacronym{mcu}{MCU}{Microcontroller Unit}
\newacronym{ml}{ML}{Machine Learning}
\newacronym{ook}{OOK}{On-Off Keying}
\newacronym{per}{PER}{Packet Error Rate}
\newacronym{rf}{RF}{Radio Frequency}
\newacronym{rssi}{RSSI}{Received Signal Strength Indicator}
\newacronym{rtl}{RTL}{Register-Transfer Level}
\newacronym{snr}{SNR}{Signal-to-Noise Ratio}
\newacronym{spi}{SPI}{Serial Peripheral Interface}
\newacronym{uart}{UART}{Universal Asynchronous Receiver/Transmitter}
\newacronym{vhdl}{VHDL}{VHSIC Hardware Description Language}

% Convenience wrappers so you can type \IDE, \RF, etc.
\providecommand{\ASK}{\gls{ask}}
\providecommand{\BGA}{\gls{bga}}
\providecommand{\BOM}{\gls{bom}}
\providecommand{\BRAM}{\gls{bram}}
\providecommand{\CRC}{\gls{crc}}
\providecommand{\CSS}{\gls{css}}
\providecommand{\DMA}{\gls{dma}}
\providecommand{\DSP}{\gls{dsp}}
\providecommand{\EDA}{\gls{eda}}
\providecommand{\FIFO}{\gls{fifo}}
\providecommand{\FPGA}{\gls{fpga}}
\providecommand{\FSK}{\gls{fsk}}
\providecommand{\FSM}{\gls{fsm}}
\providecommand{\GPIO}{\gls{gpio}}
\providecommand{\HAL}{\gls{hal}}
\providecommand{\IIC}{\gls{i2c}}
\providecommand{\IDE}{\gls{ide}}
\providecommand{\ILA}{\gls{ila}}
\providecommand{\ISM}{\gls{ism}}
\providecommand{\ISR}{\gls{isr}}
\providecommand{\LUT}{\gls{lut}}
\providecommand{\MCU}{\gls{mcu}}
\providecommand{\ML}{\gls{ml}}
\providecommand{\OOK}{\gls{ook}}
\providecommand{\PER}{\gls{per}}
\providecommand{\RF}{\gls{rf}}
\providecommand{\RSSI}{\gls{rssi}}
\providecommand{\RTL}{\gls{rtl}}
\providecommand{\SNR}{\gls{snr}}
\providecommand{\SPI}{\gls{spi}}
\providecommand{\UART}{\gls{uart}}
\providecommand{\VHDL}{\gls{vhdl}}


% Call \printhandoutglossary just before \end{document} in each handout.
\newcommand{\printhandoutglossary}{%
	\clearpage
	\printnoidxglossary[type=\acronymtype,title={Glossary}]%
}

% ── Colors (match Module 1 lesson plan) ──────────────────────
\definecolor{stevens}{RGB}{163,31,52}
\definecolor{darkgray}{RGB}{60,60,60}
\definecolor{lightgray}{RGB}{240,240,240}
\definecolor{accentblue}{RGB}{30,80,140}

\hypersetup{
	colorlinks=true,
	linkcolor=stevens,
	urlcolor=accentblue,
	citecolor=darkgray,
}

% ── Defaults for header variables (handout overrides these) ──
\providecommand{\labid}{?}
\providecommand{\labtitleshort}{Lab Handout}

% ── Running header / footer ──────────────────────────────────
\pagestyle{fancy}
\fancyhf{}
\fancyhead[L]{\small\textcolor{darkgray}{Module 1 --- Lab Handout \labid}}
\fancyhead[R]{\small\textcolor{darkgray}{\labtitleshort}}
\fancyfoot[C]{\small\thepage}
\renewcommand{\headrulewidth}{0.4pt}

% ── Section formatting ───────────────────────────────────────
\titleformat{\section}{\Large\bfseries\color{stevens}}{\thesection}{0.6em}{}
\titleformat{\subsection}{\large\bfseries\color{darkgray}}{\thesubsection}{0.6em}{}

% ── Box styles ───────────────────────────────────────────────
\tcbset{
	infobox/.style={
		enhanced, colback=blue!3, colframe=accentblue, fonttitle=\bfseries,
		title={#1}, coltitle=white, breakable, left=6pt, right=6pt, top=4pt, bottom=4pt
	},
	warning/.style={
		enhanced, colback=red!3, colframe=stevens, fonttitle=\bfseries,
		title={#1}, coltitle=white, breakable, left=6pt, right=6pt, top=4pt, bottom=4pt
	},
	deliverable/.style={
		enhanced, colback=gray!5, colframe=darkgray, fonttitle=\bfseries,
		title={#1}, coltitle=white, breakable, left=6pt, right=6pt, top=4pt, bottom=4pt
	},
}

% ── Figure helper: always draws a placeholder box (images suppressed) ─
% Usage: \labfigure{filename}{width}{caption}{label}
\newcommand{\labfigure}[4]{%
	\begin{figure}[h!]
		\centering
		\fbox{\begin{minipage}[c][2.2in][c]{#2}
			\centering\color{darkgray}\small
				\textbf{[Figure placeholder]}\\[4pt]
				\texttt{\detokenize{#1}}\\[4pt]
				\textit{#3}
		\end{minipage}}%
		\caption{#3}
		\label{#4}
	\end{figure}%
}

% ── Title banner (call at top of document body) ──────────────
\newcommand{\labheader}[2]{%
	\begin{center}
		{\LARGE\bfseries\color{stevens} Lab Handout #1}\\[3pt]
		{\Large\bfseries #2}\\[6pt]
		{\small\color{darkgray} Capstone --- Module 1: \RF{} Hardware Foundations}
	\end{center}
	\vspace{4pt}
	{\color{darkgray}\rule{\textwidth}{0.4pt}}
	\vspace{10pt}
}


\begin{document}
\labheader{1-B}{STM32 Development Environment Setup}

\section{Objective}
Install the VS Code + STMicroelectronics extension stack, connect to the Nucleo-L476RG over its on-board ST-LINK/V2.1, and successfully build, flash, and debug a ``blinky'' firmware image. On completion you will have a fully functional STM32 toolchain ready for \SPI{} bring-up in Lab Handout~1-C.

\section{Prerequisites}
\begin{itemize}[leftmargin=2em]
	\item A Windows, macOS, or Linux development machine with $\geq$10\,GB of free disk space.
	\item Administrator rights (required for ST-LINK USB driver installation on Windows).
	\item STM32 Nucleo-L476RG board and a USB Mini-B cable.
\end{itemize}

\section{Procedure}

\subsection{Step 1 --- Install VS Code}
\begin{enumerate}[leftmargin=2em]
	\item Download VS Code from \href{https://code.visualstudio.com/}{code.visualstudio.com} and install with default options.
	\item Launch VS Code and open the Extensions panel (\texttt{Ctrl+Shift+X}).
\end{enumerate}

\subsection{Step 2 --- Install the STMicroelectronics extension stack}
\begin{enumerate}[leftmargin=2em]
	\item Search the marketplace for \texttt{STM32 VS Code Extension} (publisher: STMicroelectronics). Install it.
	\item The umbrella extension will auto-install the following sub-extensions. Verify each is present and enabled:
	\begin{itemize}[leftmargin=1.2em, topsep=2pt, itemsep=1pt]
		\item \texttt{stm32cube-ide-project-manager}
		\item \texttt{stm32cube-ide-build-cmake}
		\item \texttt{stm32cube-ide-build-analyzer}
		\item \texttt{stm32cube-ide-bundles-manager}
		\item \texttt{stm32cube-ide-clangd}
		\item \texttt{stm32cube-ide-debug-core}
		\item \texttt{stm32cube-ide-debug-stlink-gdbserver}
		\item \texttt{stm32cube-ide-registers}
	\end{itemize}
	\item Recommended supporting extensions (install manually):
	\begin{itemize}[leftmargin=1.2em, topsep=2pt, itemsep=1pt]
		\item \texttt{ms-vscode.cpptools-extension-pack}
		\item \texttt{ms-vscode.cmake-tools}
		\item \texttt{marus25.cortex-debug}
		\item \texttt{dan-c-underwood.arm}
		\item \texttt{eclipse-cdt.serial-monitor}
	\end{itemize}
\end{enumerate}

\subsection{Step 3 --- Install the ARM toolchain and \HAL{} packs}
\begin{enumerate}[leftmargin=2em]
	\item Open the STM32Cube Bundles Manager from the VS Code command palette (\texttt{Ctrl+Shift+P} $\to$ ``STM32: Bundles Manager'').
	\item Install the latest ARM GCC toolchain when prompted.
	\item Install the STM32Cube\mbox{L4} firmware pack (this contains the \HAL{} drivers for the L476RG).
	\item Wait for the downloads to finish before proceeding.
\end{enumerate}

\subsection{Step 4 --- Create an empty project targeting the Nucleo-L476RG}
\begin{enumerate}[leftmargin=2em]
	\item Command palette $\to$ ``STM32: New Project''.
	\item Project type: \textbf{Empty}. Board: \textbf{NUCLEO-L476RG}. Toolchain: \textbf{CMake}.
	\item When prompted, open the integrated STM32CubeMX peripheral configurator. Enable the \textbf{LED2} user \mbox{LED} on \texttt{PA5} as a \mbox{GPIO} output. Leave all other peripherals at defaults.
	\item Click ``Generate Code''. VS Code will create \texttt{main.c} with the \HAL{} initialization stubs.
\end{enumerate}

\subsection{Step 5 --- Write and build blinky}
\begin{enumerate}[leftmargin=2em]
	\item In the main loop of \texttt{main.c}, add:
\begin{verbatim}
HAL_GPIO_TogglePin(GPIOA, GPIO_PIN_5);
HAL_Delay(500);
\end{verbatim}
	\item Command palette $\to$ ``CMake: Build''. Confirm the build completes with zero errors.
	\item Note the ELF output path reported at the end of the build log; you will use it in Step~6.
\end{enumerate}

\subsection{Step 6 --- Connect the Nucleo and flash}
\begin{enumerate}[leftmargin=2em]
	\item Plug the Nucleo into your PC via USB Mini-B. The \texttt{LD1} ST-LINK status \mbox{LED} should turn red, then blink red/green while enumerating.
	\item On Windows, Device Manager should show ``ST-Link Debug'' and ``STLink Virtual COM Port'' after enumeration. On macOS/Linux, the device appears as \texttt{/dev/cu.usbmodemXXXX} or \texttt{/dev/ttyACMX}.
	\item Command palette $\to$ ``CMake: Debug'' or press \texttt{F5}. VS Code launches the ST-LINK GDB server, flashes the ELF, halts at \texttt{main}, and waits at the first breakpoint.
	\item Press \texttt{F5} again (Continue). The green \texttt{LD2} \mbox{LED} on the Nucleo should begin blinking at 1\,Hz.
\end{enumerate}

\subsection{Step 7 --- Verify the debug features}
\begin{enumerate}[leftmargin=2em]
	\item With the target halted, open the Peripheral Registers view (from the \texttt{stm32cube-ide-registers} extension). Expand \texttt{GPIOA} and locate the \texttt{ODR} register. Confirm that single-stepping past the \texttt{TogglePin} call flips bit 5.
	\item Set a breakpoint inside the main loop. Resume. Confirm the debugger halts on the breakpoint and that variable values inspect correctly.
	\item Open the Serial Monitor panel (\texttt{eclipse-cdt.serial-monitor}), select the STLink Virtual COM port, set 115200\,8N1. No data yet --- this just verifies the port opens without errors.
\end{enumerate}

\section{Verification Checklist}
\begin{itemize}[leftmargin=2em]
	\item[$\square$] All listed ST extensions are installed and enabled.
	\item[$\square$] ARM toolchain and STM32Cube\mbox{L4} \HAL{} pack installed via Bundles Manager.
	\item[$\square$] Empty project for Nucleo-L476RG builds with zero errors.
	\item[$\square$] ELF flashes successfully over ST-LINK; \texttt{LD2} blinks at 1\,Hz.
	\item[$\square$] A breakpoint in the main loop halts execution as expected.
	\item[$\square$] The \texttt{GPIOA->ODR} register view reflects the toggling state.
	\item[$\square$] Serial Monitor opens the STLink Virtual COM port without errors.
\end{itemize}

\begin{tcolorbox}[warning={Common Pitfalls}]
	\begin{itemize}[leftmargin=1.2em]
		\item \textbf{ST-LINK not detected on Windows:} install the STSW-LINK009 driver from \texttt{st.com}. A reboot may be required.
		\item \textbf{Firmware builds but fails to flash:} update the on-board ST-LINK firmware via the ``STM32: Update ST-LINK firmware'' command.
		\item \textbf{Wrong board selected:} the L476RG has an ARM Cortex-M4 at 80\,MHz; if the build log mentions M0 or M3, you picked a different Nucleo variant.
	\end{itemize}
\end{tcolorbox}

\begin{tcolorbox}[deliverable={Lab Handout 1-B Deliverable}]
	Submit the following:
	\begin{enumerate}[leftmargin=1.2em]
		\item Screenshot of the VS Code Extensions panel showing the installed STM32 stack.
		\item Screenshot of a successful build log (terminal output).
		\item Short video (5--10\,s) of \texttt{LD2} blinking on the physical Nucleo after flash.
		\item Screenshot of the halted debugger with a breakpoint hit in \texttt{main.c}.
	\end{enumerate}
\end{tcolorbox}

\end{document}
